\chapter{2004年湖北卷（文）}
\section{单选题}
\begin{problem}
设 \(A=\{x\mid x=\sqrt{5k+1}, k\in\N\}\)，\(B=\{x\mid x\leqslant 6, x\in\mathbf{Q}\}\)，则 \(A\cap B\) 等于（\quad）

\choices
    {\(\{1,4\}\)}
    {\(\{1,6\}\)}
    {\(\{4,6\}\)}
    {\(\{1,4,6\}\)}
\answer{D}
\end{problem}

\begin{solution}
因为 \(x=\sqrt{5k+1}\) 为有理数，且 \(5k+1\) 是整数，所以 \(x\) 必为整数. 又 \(x\geqslant0\) 且 \(x\leqslant6\)，故先考察 \(x=0\)，\(1\)，\(2\)，\(3\)，\(4\)，\(5\)，\(6\). 由 \(k=(x^2-1)/5\in\N\) 检验可得 \(x=1\)，\(4\)，\(6\)，而 \(x=0\) 对应 \(k=-\frac15\notin\N\)，从而 \(A\cap B=\{1,4,6\}\)，故选 D.
\end{solution}

\begin{problem}
已知点 \(M_1\left(6,2\right)\) 和 \(M_2\left(1,7\right)\). 直线 \(y=mx-7\) 与线段 \(M_1M_2\) 的交点 \(M\) 分有向线段 \(M_1M_2\) 的比为 \(3:2\)，则 \(m\) 的值为（\quad）

\choices
    {\(-\frac{3}{2}\)}
    {\(-\frac{2}{3}\)}
    {\(\frac{1}{4}\)}
    {4}
\answer{D}
\end{problem}

\begin{solution}
由内分点公式，点 \(M\) 的坐标为
\(M=\dfrac{2M_1+3M_2}{5}=\left(\dfrac{2\times6+3\times1}{5},\dfrac{2\times2+3\times7}{5}\right)=\left(3,5\right)\).
因为 \(M\) 在直线 \(y=mx-7\) 上，所以 \(5=3m-7\)，解得 \(m=4\)，故选 D.
\end{solution}

\begin{problem}
已知函数 \(f(x)\) 在 \(x=1\) 处的导数为 3，则 \(f(x)\) 的解析式可能为（\quad）

\choices
    {\(f(x)=\left(x-1\right)^2+3\left(x-1\right)\)}
    {\(f(x)=2\left(x-1\right)\)}
    {\(f(x)=2\left(x-1\right)^2\)}
    {\(f(x)=x-1\)}
\answer{A}
\end{problem}

\begin{solution}
分别求四个选项中函数在 \(x=1\) 处的导数. 选项 A 的导数为 \(2\left(x-1\right)+3\)，在 \(x=1\) 时等于 \(3\)；选项 B 的导数恒为 \(2\)；选项 C 的导数为 \(4\left(x-1\right)\)，在 \(x=1\) 时等于 \(0\)；选项 D 的导数恒为 \(1\). 因此只有选项 A 符合，故选 A.
\end{solution}

\begin{problem}
两个圆 \(C_1: x^2+y^2+2x+2y-2=0\) 与 \(C_2: x^2+y^2-4x-2y+1=0\) 的公切线有且仅有（\quad）

\choices
    {\(1\) 条}
    {\(2\) 条}
    {\(3\) 条}
    {\(4\) 条}
\answer{B}
\end{problem}

\begin{solution}
将两圆方程分别配方，得
\(C_1:\left(x+1\right)^2+\left(y+1\right)^2=4\)，\(C_2:\left(x-2\right)^2+\left(y-1\right)^2=4\).
所以两圆半径均为 \(2\)，圆心距为
\(d=\sqrt{(2+1)^2+(1+1)^2}=\sqrt{13}\).
由于 \(0<d<2+2\)，且两圆半径相等，故两圆相交于两个不同的点. 相交两圆只有两条公切线，即选项 B.
\end{solution}

\begin{problem}
若函数 \(f(x)=a^x+b-1\)（\(a>0\)，\(a\neq 1\)）的图象经过第二，三，四象限，则一定有（\quad）

\choices
    {\(0<a<1\)，且 \(b>0\)}
    {\(a>1\)，且 \(b>0\)}
    {\(0<a<1\)，且 \(b<0\)}
    {\(a>1\)，且 \(b<0\)}
\answer{C}
\end{problem}

\begin{solution}
函数图象为 \(y=a^x+b-1\)，且 \(f(0)=b\). 若 \(a>1\)，函数递增；要同时经过第二、三象限，需有 \(b-1<0<b\)，但此时 \(x>0\) 时 \(f(x)>f(0)=b>0\)，不可能经过第四象限. 因此必有 \(0<a<1\).

当 \(0<a<1\) 时函数递减. 要经过第四象限，需有 \(b-1<0\)，即 \(b<1\)；要在 \(x<0\) 时既取正值又取负值以经过第二、三象限，还必须有 \(f(0)=b<0\). 于是 \(0<a<1\) 且 \(b<0\)，故选 C.
\end{solution}

\begin{problem}
四面体 \(ABCD\) 四个面的重心分别为 \(E\)，\(F\)，\(G\)，\(H\)，则四面体 \(EFGH\) 的表面积与四面体 \(ABCD\) 的表面积的比值是（\quad）

\choices
    {\(\frac{1}{27}\)}
    {\(\frac{1}{16}\)}
    {\(\frac{1}{9}\)}
    {\(\frac{1}{8}\)}
\answer{C}
\end{problem}

\begin{solution}
设 \(E\)，\(F\)，\(G\)，\(H\) 分别为与顶点 \(A\)，\(B\)，\(C\)，\(D\) 对面的三个顶点所组成的面的重心，即
\(E=\dfrac{B+C+D}{3}\)，\(F=\dfrac{A+C+D}{3}\)，\(G=\dfrac{A+B+D}{3}\)，\(H=\dfrac{A+B+C}{3}\).
于是
\(\overrightarrow{EF}=-\dfrac{\overrightarrow{AB}}{3}\)，\(\overrightarrow{EG}=-\dfrac{\overrightarrow{AC}}{3}\)，\(\overrightarrow{EH}=-\dfrac{\overrightarrow{AD}}{3}\).
因此四面体 \(EFGH\) 与四面体 \(ABCD\) 相似，且相似比为 \(1:3\). 表面积比等于相似比的平方，所以
\(\dfrac{S_{EFGH}}{S_{ABCD}}=\left(\dfrac{1}{3}\right)^2=\dfrac{1}{9}\)，故选 C.
\end{solution}

\begin{problem}
已知 \(\bs{a}\)，\(\bs{b}\)，\(\bs{c}\) 为非零的平面向量. 甲：\(\bs{a}\cdot\bs{b}=\bs{a}\cdot\bs{c}\)，乙：\(\bs{b}=\bs{c}\)，则（\quad）

\choices
    {甲是乙的充分条件但不是必要条件}
    {甲是乙的必要条件但不是充分条件}
    {甲是乙的充要条件}
    {甲既不是乙的充分条件也不是乙的必要条件}
\answer{B}
\end{problem}

\begin{solution}
由甲得 \(\bs{a}\cdot(\bs{b}-\bs{c})=0\). 若乙成立，即 \(\bs{b}=\bs{c}\)，则甲必成立，所以甲是乙的必要条件.

但甲不一定推出乙. 例如取 \(\bs{a}=\left(1,0\right)\)，\(\bs{b}=\left(0,1\right)\)，\(\bs{c}=\left(0,2\right)\)，则三个向量均非零，且 \(\bs{a}\cdot\bs{b}=\bs{a}\cdot\bs{c}=0\)，但 \(\bs{b}\neq\bs{c}\). 因此甲不是乙的充分条件，故选 B.
\end{solution}

\begin{problem}
已知 \(x\geqslant \frac{5}{2}\)，则 \(f(x)=\frac{x^2-4x+5}{2x-4}\) 有（\quad）

\choices
    {最大值 \(\frac{5}{4}\)}
    {最小值 \(\frac{5}{4}\)}
    {最大值 1}
    {最小值 1}
\answer{D}
\end{problem}

\begin{solution}
令 \(t=x-2\)，则 \(t\geqslant\dfrac{1}{2}>0\)，且
\(f(x)=\dfrac{\left(x-2\right)^2+1}{2\left(x-2\right)}=\dfrac{1}{2}\left(t+\dfrac{1}{t}\right)\).
由 \(t>0\)，有 \(t+\dfrac{1}{t}\geqslant2\)，等号当且仅当 \(t=1\). 所以 \(f(x)\geqslant1\)，且在 \(x=3\) 时取到最小值 \(1\). 当 \(t\to+\infty\) 时 \(f(x)\to+\infty\)，故没有最大值，选 D.
\end{solution}

\begin{problem}
设数列 \(\{a_n\}\) 的前 \(n\) 项和 \(S_n=a\left[2-\left(\frac{1}{2}\right)^{n-1}\right]-b\left[2-\left(n+1\right)\left(\frac{1}{2}\right)^{n-1}\right]\)（\(n=1\)，\(2\)，\(\cdots\)），其中 \(a\)，\(b\) 是非零常数. 则存在数列 \(\{x_n\}\)，\(\{y_n\}\) 使得（\quad）

\choices
    {\(a_n=x_n+y_n\)，其中 \(\{x_n\}\) 为等差数列，\(\{y_n\}\) 为等比数列}
    {\(a_n=x_n+y_n\)，其中 \(\{x_n\}\) 和 \(\{y_n\}\) 都为等差数列}
    {\(a_n=x_n\cdot y_n\)，其中 \(\{x_n\}\) 为等差数列，\(\{y_n\}\) 为等比数列}
    {\(a_n=x_n\cdot y_n\)，其中 \(\{x_n\}\) 和 \(\{y_n\}\) 都为等比数列}
\answer{C}
\end{problem}

\begin{solution}
令 \(r=\dfrac{1}{2}\). 由题给前 \(n\) 项和，先化为
\(S_n=2(a-b)+\bigl(b(n+1)-a\bigr)r^{n-1}\).
当 \(n\geqslant2\) 时，
\(a_n=S_n-S_{n-1}=\bigl(a+b-bn\bigr)r^{n-1}\).
当 \(n=1\) 时，题给 \(S_1=a\)，而右式也等于 \(\left(a+b-b\right)r^0=a\)，所以对所有正整数 \(n\) 都有
\(a_n=(a+b-bn)\left(\dfrac12\right)^{n-1}\).
其中 \(a+b-bn\) 是等差数列，\(\left(\dfrac12\right)^{n-1}\) 是等比数列，故可令
\(x_n=a+b-bn\)，\(y_n=\left(\dfrac12\right)^{n-1}\)，使 \(a_n=x_ny_n\). 因此选 C.
\end{solution}

\begin{problem}
若 \(1<\frac{1}{a}<\frac{1}{b}\)，则下列结论中不正确的是（\quad）

\choices
    {\(\log_a b>\log_b a\)}
    {\(\left|\log_a b+\log_b a\right|>2\)}
    {\(\left(\log_b a\right)^2<1\)}
    {\(\left|\log_a b\right|+\left|\log_b a\right|>\left|\log_a b+\log_b a\right|\)}
\answer{D}
\end{problem}

\begin{solution}
由 \(1<\dfrac1a\) 得 \(0<a<1\)，由 \(\dfrac1a<\dfrac1b\) 得 \(b<a\)，所以 \(0<b<a<1\). 令 \(u=\log_a b\)，由于 \(\ln b<\ln a<0\)，有 \(u>1\)，且
\(\log_b a=\dfrac1u\in(0,1)\).

于是 \(u>\dfrac1u\)，故选项 A 正确；由 \(u>0\) 且 \(u\neq1\)，有 \(u+\dfrac1u>2\)，故选项 B 正确；又 \(\left(\log_b a\right)^2=\dfrac1{u^2}<1\)，故选项 C 正确. 最后
\(\left|\log_a b\right|+\left|\log_b a\right|=u+\dfrac1u=\left|\log_a b+\log_b a\right|\)，不可能大于右端，故不正确的是 D.
\end{solution}

\begin{problem}
将标号为 \(1\)，\(2\)，\(\cdots\)，\(10\) 的 \(10\) 个球放入标号为 \(1\)，\(2\)，\(\cdots\)，\(10\) 的 \(10\) 个盒子里，每个盒内放一个球，恰好 \(3\) 个球的标号与其在盒子中的标号不一致的放入方法种数为（\quad）

\choices
    {120}
    {240}
    {360}
    {720}
\answer{B}
\end{problem}

\begin{solution}
恰有 \(3\) 个球的标号与盒子标号不一致，意味着先选出这 \(3\) 个盒子，共有 \(\mathrm C_{10}^3\) 种选法. 选定后，这 \(3\) 个球在这 \(3\) 个盒子中必须形成无固定点的排列；三个元素的无固定点排列只能是两个三循环，共有 \(2\) 种.
因此放入方法数为
\(\mathrm C_{10}^3\times2=120\times2=240\)，故选 B.
\end{solution}

\begin{problem}
设 \(y=f(t)\) 是某港口水的深度 \(y\)（米）关于时间 \(t\)（时）的函数，其中 \(0\leqslant t\leqslant 24\)，下表是该港口某一天从 0 时至 24 时记录的时间 \(t\) 与水深 \(y\) 的关系：

\begin{center}
\begin{tabular}{|c|c|c|c|c|c|c|c|c|c|}
\hline
\(t\) & 0 & 3 & 6 & 9 & 12 & 15 & 18 & 21 & 24 \\
\hline
\(y\) & 12 & 15.1 & 12.1 & 9.1 & 11.9 & 14.9 & 11.9 & 8.9 & 12.1 \\
\hline
\end{tabular}
\end{center}

经长期观察，函数 \(y=f(t)\) 的图象可以近似地看成函数 \(y=k+A\sin\left(\omega t+\varphi\right)\) 的图象. 下面的函数中，最能近似表示表中数据间对应关系的函数是（\quad）

\choices
    {\(y=12+3\sin\dfrac{\pi}{6}t\)，\(t\in\left[0,24\right]\)}
    {\(y=12+3\sin\left(\dfrac{\pi}{6}t+\pi\right)\)，\(t\in\left[0,24\right]\)}
    {\(y=12+3\sin\dfrac{\pi}{12}t\)，\(t\in\left[0,24\right]\)}
    {\(y=12+3\sin\left(\dfrac{\pi}{12}t+\dfrac{\pi}{2}\right)\)，\(t\in\left[0,24\right]\)}
\answer{A}
\end{problem}

\begin{solution}
由表中数据可知水深的平均值约为 \(12\)，最大值约为 \(15\)，最小值约为 \(9\)，所以模型中的 \(k=12\)，振幅 \(A=3\). 水深在 \(t=3\)，\(15\) 时达到最大值，在 \(t=9\)，\(21\) 时达到最小值，周期约为 \(12\)，故
\(\omega=\dfrac{2\pi}{12}=\dfrac{\pi}{6}\).
又 \(t=0\) 时水深为 \(12\)，且随后先上升，所以可取相位 \(\varphi=0\). 因而模型为
\(y=12+3\sin\dfrac{\pi}{6}t\)，与选项 A 一致，故选 A.
\end{solution}

\section{填空题}
\begin{problem}
\(\tan 2010^\circ\) 的值为\fillinblank{}.
\answer{\(\frac{\sqrt{3}}{3}\)}
\end{problem}

\begin{solution}
因为 \(2010^\circ=11\times180^\circ+30^\circ\)，正切函数的周期为 \(180^\circ\)，所以
\(\tan2010^\circ=\tan30^\circ=\dfrac{\sqrt{3}}{3}\).
\end{solution}

\begin{problem}
已知 \(\left(x^{\frac{1}{2}}+x^{-\frac{1}{2}}\right)^n\) 的展开式中各项系数的和是 128，则展开式中 \(x^5\) 的系数是\fillinblank{}. （以数字作答）
\answer{0}
\end{problem}

\begin{solution}
令 \(x=1\)，得到各项系数的和为 \(2^n=128=2^7\)，所以 \(n=7\). 二项展开式的通项为 \(\mathrm C_7^r\left(x^{\frac{1}{2}}\right)^{7-r}\left(x^{-\frac{1}{2}}\right)^r=\mathrm C_7^r x^{\frac{7-2r}{2}}\)，其中 \(r\) 为 \(0\) 到 \(7\) 的整数. 若其中有 \(x^5\) 项，则应有 \(\frac{7-2r}{2}=5\)，即 \(r=-\frac{3}{2}\)，这不可能是整数. 因此展开式中没有 \(x^5\) 项，其系数为 \(0\).
\end{solution}

\begin{problem}
某校有老师 \(200\) 人，男学生 \(1200\) 人，女学生 \(1000\) 人. 现用分层抽样的方法从所有师生中抽取一个容量为 \(n\) 的样本；已知从女学生中抽取的人数为 \(80\) 人，则 \(n=\)\fillinblank{}.
\answer{192}
\end{problem}

\begin{solution}
全校师生总人数为 \(200+1200+1000=2400\)，分层抽样时各层抽取人数与该层人数成比例. 设样本容量为 \(n\)，则女学生层满足
\(\dfrac{80}{n}=\dfrac{1000}{2400}=\dfrac{5}{12}\).
因此 \(n=80\times\dfrac{12}{5}=192\).
\end{solution}

\begin{problem}
设 \(A\)，\(B\) 为两个集合. 下列四个命题：

\begin{circlelist}
\item \(A\not\subseteq B\Leftrightarrow\) 对任意 \(x\in A\)，有 \(x\notin B\)；
\item \(A\not\subseteq B\Leftrightarrow A\cap B=\varnothing\)；
\item \(A\not\subseteq B\Leftrightarrow A\not\supseteq B\)；
\item \(A\not\subseteq B\Leftrightarrow\) 存在 \(x\in A\)，使得 \(x\notin B\).
\end{circlelist}

其中真命题的序号是\fillinblank{}. （把符合要求的命题序号都填上）
\answer{4}
\end{problem}

\begin{solution}
命题 \(\circled{1}\) 把“\(A\not\subseteq B\)”错误地表述成对任意 \(x\in A\) 都有 \(x\notin B\). 例如取 \(A=\{1,2\}\)，\(B=\{2\}\)，则 \(A\not\subseteq B\) 为真，但 \(2\in A\cap B\)，所以命题 \(\circled{1}\) 不是真命题. \(A\not\subseteq B\) 的正确含义只是存在至少一个 \(x\in A\) 不属于 \(B\).

命题 \(\circled{2}\) 也不正确. 例如 \(A=\{1,2\}\)，\(B=\{2\}\) 时，\(A\not\subseteq B\)，但 \(A\cap B=\{2\}\neq\varnothing\). 命题 \(\circled{3}\) 不正确，例如 \(A=\{1\}\)，\(B=\{1,2\}\) 时，\(A\not\subseteq B\) 为假，而 \(A\not\supseteq B\) 为真.

命题 \(\circled{4}\) 正是集合包含关系否定的定义：\(A\not\subseteq B\) 当且仅当存在 \(x\in A\) 使 \(x\notin B\). 因此真命题的序号只有 \(4\).
\end{solution}

\begin{problem}
已知 \(6\sin^2\alpha+\sin\alpha\cos\alpha-2\cos^2\alpha=0\)，\(\alpha\in\left[\frac{\pi}{2},\pi\right]\)，则 \(\sin\left(2\alpha+\frac{\pi}{3}\right)=\)\fillinblank{}.
\answer{\(\frac{5\sqrt{3}-12}{26}\)}
\end{problem}

\begin{solution}
当 \(\alpha=\frac{\pi}{2}\) 时，原方程左边为 \(6\)，不成立，故 \(\cos\alpha\neq0\). 原方程同除以 \(\cos^2\alpha\)，得 \(6\tan^2\alpha+\tan\alpha-2=0\)，即 \((2\tan\alpha-1)(3\tan\alpha+2)=0\). 因为 \(\alpha\in\left[\frac{\pi}{2},\pi\right]\)，所以 \(\tan\alpha\leqslant0\)，从而 \(\tan\alpha=-\frac{2}{3}\). 又 \(\sin\alpha\geqslant0\)，\(\cos\alpha\leqslant0\)，故 \(\sin\alpha=\frac{2\sqrt{13}}{13}\)，\(\cos\alpha=-\frac{3\sqrt{13}}{13}\). 于是 \(\sin2\alpha=-\frac{12}{13}\)，\(\cos2\alpha=\frac{5}{13}\)，所以 \(\sin\left(2\alpha+\frac{\pi}{3}\right)=\sin2\alpha\cos\frac{\pi}{3}+\cos2\alpha\sin\frac{\pi}{3}=\frac{5\sqrt{3}-12}{26}\).
\end{solution}

\section{解答题}
\begin{problem}
如图，在棱长为 \(1\) 的正方体 \(ABCD-A_1B_1C_1D_1\) 中，\(AC\) 与 \(BD\) 交于点 \(E\)，\(C_1B\) 与 \(CB_1\) 交于点 \(F\).

\begin{enumerate}
\item 求证：\(A_1C\perp\) 平面 \(BDC_1\)；

\item 求二面角 \(B-EF-C\) 的大小. （结果用反三角函数值表示）
\end{enumerate}

\bitmapfigure[max width=0.44\linewidth]{img/2004/hubei_liberal/q18_fig1.png}
\end{problem}

\begin{solution}
建立空间直角坐标系，使 \(A\left(0,0,0\right)\)，\(B\left(1,0,0\right)\)，\(C\left(1,1,0\right)\)，\(D\left(0,1,0\right)\)，\(A_1\left(0,0,1\right)\)，\(B_1\left(1,0,1\right)\)，\(C_1\left(1,1,1\right)\).

\begin{enumerate}
\item 有 \(\overrightarrow{A_1C}=\left(1,1,-1\right)\)，\(\overrightarrow{DB}=\left(1,-1,0\right)\)，\(\overrightarrow{DC_1}=\left(1,0,1\right)\). 因此 \(\overrightarrow{A_1C}\cdot\overrightarrow{DB}=0\)，\(\overrightarrow{A_1C}\cdot\overrightarrow{DC_1}=0\). 直线 \(DB\) 与 \(DC_1\) 是平面 \(BDC_1\) 内相交的两条直线. 又直线 \(A_1C\) 与平面 \(BDC_1\) 相交于点 \(P\left(\frac{2}{3},\frac{2}{3},\frac{1}{3}\right)\)，因为 \(P=A_1+\frac{2}{3}\overrightarrow{A_1C}\)，且 \(\overrightarrow{DP}=\frac{1}{3}\overrightarrow{DB}+\frac{1}{3}\overrightarrow{DC_1}\). 所以 \(A_1C\perp\) 平面 \(BDC_1\).
\item 由对角线的中点公式，\(E\left(\frac{1}{2},\frac{1}{2},0\right)\). 直线 \(C_1B\) 与 \(CB_1\) 的参数方程分别为 \(\left(1,1-t,1-t\right)\) 和 \(\left(1,1-s,s\right)\)，联立得 \(t=s=\frac{1}{2}\)，所以 \(F\left(1,\frac{1}{2},\frac{1}{2}\right)\). 设 \(H\) 为 \(EF\) 的中点，则 \(H\left(\frac{3}{4},\frac{1}{2},\frac{1}{4}\right)\). 因为 \(\overrightarrow{EF}=\left(\frac{1}{2},0,\frac{1}{2}\right)\)，\(\overrightarrow{HB}=\left(\frac{1}{4},-\frac{1}{2},-\frac{1}{4}\right)\)，\(\overrightarrow{HC}=\left(\frac{1}{4},\frac{1}{2},-\frac{1}{4}\right)\)，所以 \(\overrightarrow{HB}\cdot\overrightarrow{EF}=0\)，\(\overrightarrow{HC}\cdot\overrightarrow{EF}=0\). 因此 \(\angle BHC\) 是二面角 \(B-EF-C\) 的平面角. 又 \(\overrightarrow{HB}\cdot\overrightarrow{HC}=-\frac{1}{8}\)，\(|\overrightarrow{HB}|^2=|\overrightarrow{HC}|^2=\frac{3}{8}\)，从而 \(\cos\angle BHC=-\frac{1}{3}\). 所以二面角 \(B-EF-C\) 的大小为 \(\arccos\left(-\frac{1}{3}\right)\).
\end{enumerate}
\end{solution}

\begin{problem}
如图，在 \(\mathrm{Rt}\triangle ABC\) 中，已知 \(BC=a\)，若长为 \(2a\) 的线段 \(PQ\) 以点 \(A\) 为中点，问 \(\overrightarrow{PQ}\) 与 \(\overrightarrow{BC}\) 的夹角 \(\theta\) 取何值时 \(\overrightarrow{BP}\cdot\overrightarrow{CQ}\) 的值最大？ 并求出这个最大值.

\bitmapfigure[max width=0.30\linewidth]{img/2004/hubei_liberal/q19_fig1.png}
\end{problem}

\begin{solution}
设 \(\bs{u}=\overrightarrow{AB}\)，\(\bs{v}=\overrightarrow{AC}\)，\(\bs{p}=\overrightarrow{AP}\). 因为 \(A\) 是 \(PQ\) 的中点，所以 \(\overrightarrow{AQ}=-\bs{p}\)，\(\overrightarrow{PQ}=-2\bs{p}\). 又 \(|PQ|=2a\)，故 \(|\bs{p}|=a\). 在直角三角形 \(ABC\) 中，\(\bs{u}\cdot\bs{v}=0\)，且 \(|\bs{v}-\bs{u}|=|\overrightarrow{BC}|=a\).

有 \(\overrightarrow{BP}=\bs{p}-\bs{u}\)，\(\overrightarrow{CQ}=-\bs{p}-\bs{v}\)，所以 \(\overrightarrow{BP}\cdot\overrightarrow{CQ}=(\bs{p}-\bs{u})\cdot(-\bs{p}-\bs{v})=-|\bs{p}|^2+\bs{p}\cdot(\bs{u}-\bs{v})=-a^2+\bs{p}\cdot(\bs{u}-\bs{v})\). 另一方面，\(\overrightarrow{PQ}\cdot\overrightarrow{BC}=(-2\bs{p})\cdot(\bs{v}-\bs{u})=2\bs{p}\cdot(\bs{u}-\bs{v})\)，故 \(\bs{p}\cdot(\bs{u}-\bs{v})=\frac{1}{2}|PQ||BC|\cos\theta=a^2\cos\theta\). 于是 \(\overrightarrow{BP}\cdot\overrightarrow{CQ}=a^2(\cos\theta-1)\leqslant0\). 当且仅当 \(\cos\theta=1\)，即 \(\theta=0\) 时取等号. 因此 \(\theta=0\) 时该数量积取得最大值，最大值为 \(0\).
\end{solution}

\begin{problem}
直线 \(l: y=kx+1\) 与双曲线 \(C: 2x^2-y^2=1\) 的右支交于不同的两点 \(A\)，\(B\).

\begin{enumerate}
\item 求实数 \(k\) 的取值范围；

\item 是否存在实数 \(k\)，使得以线段 \(AB\) 为直径的圆经过双曲线 \(C\) 的右焦点 \(F\)？ 若存在，求出 \(k\) 的值. 若不存在，说明理由.
\end{enumerate}
\end{problem}

\begin{solution}
设直线与双曲线的两个交点的横坐标为 \(x_1\)，\(x_2\).

\begin{enumerate}
\item 将 \(y=kx+1\) 代入 \(2x^2-y^2=1\)，得 \((2-k^2)x^2-2kx-2=0\). 要有两个不同的实交点，必须有判别式 \(\Delta=4(4-k^2)>0\)，所以 \(|k|<2\). 当 \(k=\pm\sqrt{2}\) 时，二次项系数为 \(0\)，方程退化为一次方程，只能得到一个交点，故还需排除这两个值. 又因为两个交点在双曲线的右支上，所以 \(x_1>0\)，\(x_2>0\). 当 \(|k|<\sqrt{2}\) 时，\(x_1x_2=-\frac{2}{2-k^2}<0\)，不可能两个横坐标都为正；当 \(\sqrt{2}<|k|<2\) 时，\(x_1x_2>0\)，且 \(x_1+x_2=\frac{2k}{2-k^2}\). 因分母为负，要使两根均为正，必须有 \(k<0\). 反之，\(-2<k<-\sqrt{2}\) 时判别式为正、两根之积为正且两根之和为正，故两根均为正. 因此 \(k\) 的取值范围为 \((-2,-\sqrt{2})\).
\item 双曲线可写为 \(\frac{x^2}{1/2}-y^2=1\)，所以其右焦点为 \(F\left(\frac{\sqrt{6}}{2},0\right)\). 记 \(f=\frac{\sqrt{6}}{2}\)，则由韦达定理，\(x_1+x_2=\frac{2k}{2-k^2}\)，\(x_1x_2=-\frac{2}{2-k^2}\). 以 \(AB\) 为直径的圆经过 \(F\)，等价于 \(\overrightarrow{FA}\cdot\overrightarrow{FB}=0\)，即
\(\left(f-x_1\right)\left(f-x_2\right)+\left(kx_1+1\right)\left(kx_2+1\right)=0\).
代入两根的和与积，得 \(f^2+1+(1+k^2)x_1x_2+(k-f)(x_1+x_2)=0\). 再代入 \(f^2=\frac{3}{2}\)，化简为 \(5k^2+2\sqrt{6}k-6=0\)，所以 \(k=\frac{6-\sqrt{6}}{5}\) 或 \(k=-\frac{6+\sqrt{6}}{5}\). 结合第（1）问的范围，只有 \(k=-\frac{6+\sqrt{6}}{5}\) 符合. 因此存在这样的实数 \(k\)，且其值为 \(-\frac{6+\sqrt{6}}{5}\).
\end{enumerate}
\end{solution}

\begin{problem}
为防止某突发事件发生，有甲，乙，丙，丁四种相互独立的预防措施可供采用，单独采用甲，乙，丙，丁预防措施后此突发事件不发生的概率（记为 \(P\)）和所需费用如下表：

\begin{center}
\begin{tabular}{|c|c|c|c|c|}
\hline
预防措施 & 甲 & 乙 & 丙 & 丁 \\
\hline
\(P\) & 0.9 & 0.8 & 0.7 & 0.6 \\
\hline
费用（万元） & 90 & 60 & 30 & 10 \\
\hline
\end{tabular}
\end{center}

预防方案可单独采用一种预防措施或联合采用几种预防措施. 在总费用不超过 120 万元的前提下，请确定一个预防方案，使得此突发事件不发生的概率最大.
\end{problem}

\begin{solution}
若联合采用若干措施，由于四种措施相互独立，联合方案使事件不发生的概率为 \(1\) 减去所有措施均失效的概率. 甲、乙、丙、丁单独失效的概率分别为 \(0.1\)，\(0.2\)，\(0.3\)，\(0.4\).

在费用不超过 \(120\) 万元的方案中，单独采用甲、乙、丙、丁的概率分别为 \(0.9\)，\(0.8\)，\(0.7\)，\(0.6\)；两项联合方案中，甲丙、甲丁、乙丙、乙丁、丙丁的费用和概率分别为 \(120\) 万元和 \(0.97\)，\(100\) 万元和 \(0.96\)，\(90\) 万元和 \(0.94\)，\(70\) 万元和 \(0.92\)，\(40\) 万元和 \(0.88\). 三项联合方案中只有乙丙丁的费用不超过预算，其费用为 \(100\) 万元，概率为 \(1-0.2\times0.3\times0.4=0.976\)；含甲的三项方案费用分别为 \(130\)，\(160\)，\(180\) 万元，均超过预算. 四项联合费用为 \(190\) 万元，也超过预算. 因此最大概率为 \(0.976\)，应联合采用乙、丙、丁三种预防措施.
\end{solution}

\begin{problem}
已知 \(b>-1\)，\(c>0\)，函数 \(f(x)=x+b\) 的图象与函数 \(g(x)=x^2+bx+c\) 的图象相切.

\begin{enumerate}
\item 求 \(b\) 与 \(c\) 的关系式（用 \(c\) 表示 \(b\)）；

\item 设函数 \(F(x)=f(x)g(x)\) 在 \(\left(-\infty,+\infty\right)\) 内有极值点，求 \(c\) 的取值范围.
\end{enumerate}
\end{problem}

\begin{solution}
\begin{enumerate}
\item 函数 \(f(x)\) 与 \(g(x)\) 的图象相切，等价于方程 \(x+b=x^2+bx+c\) 有重根，即方程 \(x^2+(b-1)x+c-b=0\) 的判别式为零. 因此 \((b-1)^2-4(c-b)=0\)，从而 \(c=\frac{(b+1)^2}{4}\). 由于 \(b>-1\)，所以 \(b+1>0\)，也可写成 \(b=2\sqrt{c}-1\).
\item 展开得 \(F(x)=x^3+2bx^2+(b^2+c)x+bc\)，所以 \(F'(x)=3x^2+4bx+b^2+c\). 这是开口向上的二次函数. 要使三次函数 \(F(x)\) 在整个实数范围内有极值点，必须且只需 \(F'(x)\) 有两个不相等的实根：此时 \(F'\) 在两个根处变号，\(F\) 分别取得极大值和极小值；若根不为两个不相等的实根，则 \(F'\) 不发生两次变号，\(F\) 没有极值点. 因此
\(\Delta=(4b)^2-12(b^2+c)=4(b^2-3c)>0\).
代入 \(c=\frac{(b+1)^2}{4}\)，得 \(b^2-6b-3>0\)，故 \(b<3-2\sqrt{3}\) 或 \(b>3+2\sqrt{3}\). 结合 \(b>-1\)，得到 \(-1<b<3-2\sqrt{3}\) 或 \(b>3+2\sqrt{3}\). 又因 \(c=\frac{(b+1)^2}{4}\)，且 \(b+1>0\)，所以
\(0<c<(2-\sqrt{3})^2=7-4\sqrt{3}\) 或 \(c>(2+\sqrt{3})^2=7+4\sqrt{3}\).
因此 \(c\) 的取值范围为 \(\left(0,7-4\sqrt{3}\right)\cup\left(7+4\sqrt{3},+\infty\right)\).
\end{enumerate}
\end{solution}
